\centerline{\bf \Huge Принцип Дирихле I}

\begin{flushright}
        \scriptsize{ДИРИХЛЕЕЕЕЕЕЕЕЕЕ! - Алексей Владимирович Савватеев}
\end{flushright}


\problem{\textbf{1}} В школе 45 классов и 2025 учащихся. Докажите, что есть класс, в котором не менее 45 учеников.

\begin{flushright} \textbf{1 балл}\\ 
\end{flushright}


\problem{\textbf{2}} В квадратном ковре со стороной 45 метров моль проела 2024 дырки. Докажите, что из этого ковра можно вырезать коврик со стороной 1 метр, в котором дырок не будет.

\begin{flushright} \textbf{1 балл}\\ 
\end{flushright}

\problem{\textbf{3}} Плоскость раскрашена в два цвета. Докажите, что найдутся 2 точки на расстоянии 2025 сантиметров друг от друга, раскрашенные в один цвет.

\begin{flushright} \textbf{2 балла}\\ 
\end{flushright}

\problem{\textbf{4}} Верно ли, что в нашей аудитории есть по крайне мере два человека, имеющие одинаковое число друзей в этой аудитории? Верно ли это для аудитории 10 класса?

\begin{flushright} \textbf{2 балла}\\ 
\end{flushright}

\problem{\textbf{5}} Каждая клетка таблицы 2026×2026 покрашена в один из 2025 цветов. За ход можно взять строку или столбец и, если там есть две клетки одного цвета, перекрасить эту строку или столбец в этот цвет. Всегда ли можно за несколько ходов покрасить всю таблицу в один цвет?

\begin{flushright} \textbf{2 балла}\\ 
\end{flushright}

\problem{\textbf{6}} Можно ли расставить на окружности числа 1, 2, 3, 4, 5, 6, 7, 8, 9, 10, 11, 12 так, чтобы разность между двумя рядом стоящими числами была 3, 4 или 5?

\begin{flushright} \textbf{3 балла}\\ 
\end{flushright}